\documentclass[12pt,a4paper]{article}
\usepackage{fontspec}
\setmainfont{DejaVu Serif}
\setsansfont{DejaVu Sans}
\setmonofont{DejaVu Sans Mono}
\usepackage{geometry}
\usepackage{hyperref}
\usepackage{booktabs}
\usepackage{longtable}
\usepackage{array}
\usepackage{enumitem}
\usepackage{xcolor}
\usepackage{titlesec}
\usepackage{setspace}

\geometry{a4paper, left=24mm, right=24mm, top=28mm, bottom=28mm}
\onehalfspacing

\definecolor{keywordblue}{rgb}{0.0,0.2,0.7}
\titleformat{\section}{\large\bfseries\color{keywordblue}}{\thesection}{1em}{}
\newcolumntype{P}[1]{>{\raggedright\arraybackslash}p{#1}}

\title{CodeAlchemist -- 8. Hafta Eklenenler ve Uygulama Gunlugu}
\author{Dila KEMER}
\date{Mart 2026}

\begin{document}
\maketitle
\tableofcontents
\newpage

\section{Ama\c{c}}
Bu belge, 8. hafta calismalarinda talep edilen dort ana gelistirmenin proje kod tabanina nasil eklendigini,
hangi dosyalara dokunuldugunu ve basta belirlenen hedeften canliya alinabilir bir MVP'ye kadar olan sureci uctan uca kaydeder.

\begin{enumerate}[leftmargin=*]
  \item Baglam Optimizasyon Asistani
  \item Patch Guvenlik Seviyesi
  \item Paylasimda Canli Isbirligi (Yorum + Onay/Revizyon Dongusu)
  \item Free/Premium Plan ve Kota Yonetimi
\end{enumerate}

\section{Baslangic Analizi ve Kapsamlama}
\subsection{Incelenen Alanlar}
Ilk adimda mevcut kod tabaninda entegrasyon noktalarini bulmak icin istemci ve sunucu dosyalari tarandi.

\begin{itemize}[leftmargin=*]
  \item \texttt{client/src/components/ChatInterface.jsx}: Context bar, prompt gonderim akisi, patch uygulama/undo
  \item \texttt{client/src/App.jsx}: Ana \texttt{onAsk} akisi, collaboration paylasim akisi, arac cekmecesi
  \item \texttt{server/app.py}: \texttt{/api/ask}, \texttt{/api/blend}, collaboration endpointleri
  \item \texttt{server/models.py}: SharedSession ve ilgili veri modelleri
\end{itemize}

\subsection{MVP Siniri}
8. hafta icin secilen strateji: mevcut davranisi bozmadan, yeni ozellikleri feature-level ek katman olarak sisteme almak.
Bu sayede mevcut hafta 6/7 kazanimi korunurken yeni kabiliyetler devreye alindi.

\section{Uygulanan Degisiklikler}

\subsection{H8.1 -- Baglam Optimizasyon Asistani (Prompt Oncesi)}
\textbf{Dosya:} \texttt{client/src/components/ChatInterface.jsx}

Eklenenler:
\begin{itemize}[leftmargin=*]
  \item Prompt optimizasyonu icin yardimci fonksiyonlar:
  \begin{itemize}[leftmargin=*]
    \item tekrarli satir temizleme
    \item gereksiz bosluklari sikistirma
    \item token oncesi/sonrasi kestirimi
  \end{itemize}
  \item \texttt{Auto Optimize: ON/OFF} anahtari
  \item Yuksek baglam baskisinda gonderim oncesi onerme karti
  \item Uc aksiyon:
  \begin{itemize}[leftmargin=*]
    \item \texttt{Apply Optimized Text}
    \item \texttt{Send Optimized}
    \item \texttt{Send Original}
  \end{itemize}
\end{itemize}

Kazan\i m:
\begin{itemize}[leftmargin=*]
  \item Context bar yalnizca gosteren bir widget olmaktan cikti, aksiyon ureten bir yardimciya donustu.
  \item Kullanici, promptu gondermeden once token tasarrufunu sayisal olarak gorebiliyor.
\end{itemize}

\subsection{H8.2 -- Patch Guvenlik Seviyesi}
\textbf{Dosya:} \texttt{client/src/components/ChatInterface.jsx}

Eklenenler:
\begin{itemize}[leftmargin=*]
  \item Patch risk skoru hesabi (0--100)
  \item Skor girdileri:
  \begin{itemize}[leftmargin=*]
    \item satir farki
    \item silinme orani
    \item kritik anahtar kelime temasi (auth/token/secret/billing)
    \item test sinyali
    \item buyuk patch boyutu
  \end{itemize}
  \item Modlar:
  \begin{itemize}[leftmargin=*]
    \item \texttt{safe}
    \item \texttt{review}
    \item \texttt{block}
  \end{itemize}
  \item Patch Safety modal:
  \begin{itemize}[leftmargin=*]
    \item risk gerekceleri listesi
    \item review akisi icin kontrollu devam
    \item block durumunda uygulamayi durdurma
  \end{itemize}
\end{itemize}

Ek duzeltme:
\begin{itemize}[leftmargin=*]
  \item Model karsilastirma kutularindaki dogrudan \texttt{setCode} uygulamasi kaldirildi; tum patch akisi risk kapisindan gecirildi.
\end{itemize}

\subsection{H8.3 -- Free/Premium Plan ve Kota Yonetimi}
\textbf{Dosyalar:} \texttt{server/app.py}, \texttt{client/src/App.jsx}

Sunucu tarafi:
\begin{itemize}[leftmargin=*]
  \item Plan limitleri eklendi:
  \begin{itemize}[leftmargin=*]
    \item Free: gunluk istek + aylik token limiti
    \item Premium: yuksek limit
  \end{itemize}
  \item Kota tuketim ve kontrol yardimcilari yazildi
  \item \texttt{/api/ask} ve \texttt{/api/blend} icin kota kapisi eklendi
  \item Yeni endpointler:
  \begin{itemize}[leftmargin=*]
    \item \texttt{GET /api/billing/usage}
    \item \texttt{POST /api/billing/plan}
  \end{itemize}
\end{itemize}

Istemci tarafi:
\begin{itemize}[leftmargin=*]
  \item Plan ve kullanim bilgisi cekme akisi
  \item Tool drawer icinde plan karti
  \item Free/Premium gecis butonlari (MVP simule plan gecisi)
  \item Kota asiminda yukseltime yonlendiren hata mesaji
\end{itemize}

\subsection{H8.4 -- Canli Isbirligi: Yorum + Onay/Revizyon Dongusu}
\textbf{Dosyalar:} \texttt{server/models.py}, \texttt{server/app.py}, \texttt{client/src/App.jsx}

Veri modeli:
\begin{itemize}[leftmargin=*]
  \item \texttt{CollaborationReview}: oturum review durumu
  \item \texttt{CollaborationComment}: yorum akisi
\end{itemize}

API katmani:
\begin{itemize}[leftmargin=*]
  \item \texttt{GET /api/collaboration/session/<token>/review}
  \item \texttt{POST /api/collaboration/session/<token>/review/comment}
  \item \texttt{POST /api/collaboration/session/<token>/review/status}
\end{itemize}

UI katmani:
\begin{itemize}[leftmargin=*]
  \item Canli collaboration paneline review status eklendi
  \item Durum gecisleri: \texttt{open}, \texttt{revision\_requested}, \texttt{approved}
  \item Inline yorum yazma ve son yorumlari gosterme
\end{itemize}

\section{Surec Gunlugu (Bastan Sona)}
\begin{longtable}{@{} P{1.3cm} P{4.2cm} P{7.2cm} @{} }
\toprule
Adim & Islem & Sonuc \\
\midrule
1 & Kod tabani taramasi & Mevcut context, patch, collaboration, ask/blend akislari netlestirildi \\
2 & Veri modeli genisletme & Isbirligi inceleme/yorum tablolari eklendi \\
3 & Sunucu kota katmani & Plan bazli gunluk/aylik limit kontrolu aktif edildi \\
4 & Ask/Blend entegrasyonu & Kota asiminda kontrollu 429 + yonlendirme mesaji \\
5 & Chat optimizasyonu & Prompt oncesi optimize onerileri ve manuel/otomatik secenekler eklendi \\
6 & Patch guvenlik kapisi & Risk skoru ve review/block modal mekanizmasi eklendi \\
7 & Collaboration paneli & Canli yorum + onay/revizyon dongusu UI'a eklendi \\
8 & Dogrulama & Degisen dosyalarda statik hata kontrolu temiz cikti \\
\bottomrule
\end{longtable}

\section{Eklenen Endpoint Ozeti}
\begin{itemize}[leftmargin=*]
  \item \texttt{GET /api/billing/usage}
  \item \texttt{POST /api/billing/plan}
  \item \texttt{GET /api/collaboration/session/<token>/review}
  \item \texttt{POST /api/collaboration/session/<token>/review/comment}
  \item \texttt{POST /api/collaboration/session/<token>/review/status}
\end{itemize}

\section{Dogrulama Notu}
Degistirilen ana dosyalar icin hata taramasi yapildi ve yeni sentaks/derleme hatasi tespit edilmedi.

\section{Bilinen Sinirlar ve Sonraki Iterasyon}
\begin{enumerate}[leftmargin=*]
  \item Premium gecisi su an MVP seviyesinde (odeme entegrasyonu sonraki adim).
  \item Prompt optimizasyonu kural tabanli; ileride model destekli semantic compression eklenebilir.
  \item Patch risk skoru heuristic tabanli; dosya/AST tabanli analiz ile hassasiyet arttirilabilir.
  \item Collaboration review paneli polling ile calisiyor; sonraki iterasyonda websocket/SSE ile gercek zamanli yapilabilir.
\end{enumerate}

\section{Sonuc}
8. hafta implementasyonu ile sistem, \textit{gorunurluk} seviyesinden \textit{aksiyonlanabilir guvenlik ve maliyet kontrolu} seviyesine tasindi.
Kullanici artik promptu optimize ederek gonderebiliyor, riskli patchleri kontrollu uygulatabiliyor,
paylasilan oturumlarda yorum-onay dongusu yurutebiliyor ve plan/kota durumunu net olarak gorebiliyor.

\end{document}
